\section{Preface}
While reviewing this document, it is useful to be aware of some select syntax used to distinguish between different types of objects:

\underline{Files and folders} will appear using quotations marks \say{folder/file}. A specified folder will always end with a forward slash, while a file ends with an extension (e.g., .txt).\\[\baselineskip]
\underline{Functions and variables} used in either the context of python or bash are distinguished as highlighted \code{code}. These will always be items that can be typed into an interpreter that will produce something.\\[\baselineskip]
\underline{Links} within the text will appear \textbf{bolded}. References to tables and figures will contain the index number whereas sections will refer to the section name.\\[\baselineskip]
\underline{Warnings and notes} are spread throughout the text to alert the reader to possible pitfalls or mistakes that have at one point been made. When learning about a new section, it is best practice to carefully review the warnings.
\warn{This is an example of what a warning will look like.}


Acronyms will be defined throughout the text the first time they are used. However, a collection and summary of them are given below:
\begin{itemize}
    \item GC - Gas Chromatograph
    \item CLI - Command Line Interface
    \item MFC - Mass Flow Controller
    \item TC - Thermocouple
    \item ADS - Adsorbent Trap
\end{itemize}

\pagebreak